\documentclass[11pt, a4paper]{article}

\usepackage[utf-8]{inputenc}
\usepackage[ngerman]{babel}
\usepackage[margin=1in]{geometry}
\usepackage{setspace}

\setstretch{1.15}

\title{Klassische Simulation und experimentelle Analyse von Grovers Algorithmus}
\author{Veselin Petrov Nikolaev}
\date{März 2026}

\begin{document}

\maketitle

\section*{Zusammenfassung}

Dieses Projekt zielt darauf ab, ein tiefes theoretisches und experimentelles Verständnis von Quantenberechnung durch die klassische Simulation von Grovers Suchalgorithmus zu entwickeln. Das Projekt umfasst eine umfangreiche mathematische Analyse der Grundlagen der Quanteninformatik, einschließlich Qubits als Vektoren im Hilbert-Raum, unitäre Transformationen, Superposition, Verschränkung und Quantenmessung.

Die zentrale Komponente des Projekts ist eine tiefgreifende Untersuchung von Grovers Algorithmus. Dies beinhaltet die formale Herleitung der Oracle- und Diffusionsoperatoren, eine geometrische Interpretation der Zustandsrotation in einem zweidimensionalen Unterraum und numerisch verifizierte Beispiele unter Verwendung von Matrixformalismus. Der Fokus liegt nicht nur auf der Implementierung, sondern auf dem vollständigen Verständnis der inneren Struktur und des genauen Mechanismus der Amplitudenverstärkung.

Der Algorithmus wurde mit Qiskit Aer als vollständiger Zustandsvektorsimulator implementiert und auf klassischer Hardware simuliert. Eine systematische experimentelle Evaluierung wurde durchgeführt, um Ausführungszeit, Speicherverbrauch und Skalierbarkeit als Funktion der Anzahl von Qubits zu messen. Die Experimente umfassten die Messung des klassischen Suchalgorithmus, die Bestätigung des optimalen Iterationsverhaltens und die systematische Analyse der Schaltungstiefe.

Die Skalierungsstudien zeigten, dass der Speicherverbrauch exponentiell mit der Qubit-Anzahl wächst und das primäre Engpass für die klassische Simulation darstellt. Das Projekt evaluierte die Leistung auf verschiedenen Rechnerarchitekturen: lokale CPUs, GPUs und Hochleistungsrecheninfrastrukturen (HPC) des Universitätsnetzes, einschließlich des AION-CPU-Clusters und des IRIS-GPU-Clusters mit Tesla-V100-GPUs. GPU-Beschleunigung erreichte einen konsistenten Vorteil von etwa 12× bei größeren Qubit-Anzahlen, überwand aber nicht die grundlegende exponenzielle Speicherwand.

Das Projekt demonstriert, dass Grovers Algorithmus tatsächlich eine quadratische Beschleunigung (O(√N) gegenüber O(N)) gegenüber klassischen Suchmethoden bietet, wie durch Query-Komplexitätsmessungen bestätigt. Allerdings wird diese theoretische Überlegenheit durch praktische klassische Simulationsgrenzen erheblich gemindert. Das Projekt klärt sowohl die Stärken als auch die praktischen Grenzen von Grovers quadratischer Beschleunigung auf und zeigt, wie theoretische Quantenvorteile unter Berücksichtigung klassischer Ressourcenbeschränkungen neu bewertet werden müssen.

\end{document}
